\documentclass[11pt,a4paper]{article}
\usepackage[utf8]{inputenc}
\usepackage[T1]{fontenc}
\usepackage[french]{babel}
\usepackage{amsmath,amssymb}
\usepackage{geometry,enumitem,booktabs,xcolor}
\usepackage{tcolorbox}
\usepackage{tikz}
\usetikzlibrary{arrows.meta,positioning,calc,decorations.pathreplacing}
\usepackage{hyperref}

\geometry{margin=2cm}
\hypersetup{colorlinks=true,linkcolor=blue!60!black,urlcolor=blue!70!black}
\tcbuselibrary{skins}
\newcommand{\PsiNN}{$\Psi$-NN}
\newcommand{\R}{\mathbf{R}}
\newcommand{\W}{\mathbf{W}}

\title{\textbf{Explication schématique}\\[3pt]
\large\PsiNN{}\,: Découverte automatique de la structure des PINNs\\
via distillation de connaissances\\[2pt]
\normalsize Liu et al.\ -- \textit{Nature Communications} (2025) 16:9558}
\author{Réunion 3 -- PFE}
\date{\today}

\begin{document}
\maketitle
\tableofcontents
\newpage

%=================================================================
\section{Le problème en un schéma}
%=================================================================

\begin{center}
\begin{tikzpicture}[
    every node/.style={font=\small},
    mbox/.style={draw,rounded corners=3pt,minimum height=0.9cm,align=center,text width=3.8cm},
    arr/.style={-{Stealth[length=2.5mm]},thick},
    lbl/.style={font=\scriptsize,midway,above}
]
\node[mbox,fill=red!8] (pinn)  {PINN\\physique \textbf{uniquement dans la perte}};
\node[mbox,fill=orange!10,right=1.2cm of pinn] (post) {PINN-post\\symétrie manuelle\\en couche de sortie};
\node[mbox,fill=green!10,right=1.2cm of post] (psi)  {\PsiNN{}\\physique \textbf{dans}\\les poids du réseau};

\draw[arr,red!60]   (pinn) -- node[lbl]{\color{red!60}limité} (post);
\draw[arr,green!60!black] (post) -- node[lbl]{\color{green!60!black}automatisé} (psi);

\node[below=0.3cm of pinn,font=\scriptsize,text width=3.8cm,align=center,gray] {MLP générique. Aucune conscience structurelle. La régularisation entre en conflit avec la perte EDP.};
\node[below=0.3cm of post,font=\scriptsize,text width=3.8cm,align=center,gray] {Ajoute des contraintes dures après la sortie. Nécessite une expertise. Seulement partiel.};
\node[below=0.3cm of psi,font=\scriptsize,text width=3.8cm,align=center,gray]  {Découvre la structure automatiquement. L'intègre dans les poids via $\R$. Transférable.};
\end{tikzpicture}
\end{center}

%=================================================================
\section{Le pipeline \PsiNN{} (3 étapes)}
%=================================================================

\begin{figure}[h!]
\centering
\begin{tikzpicture}[
    stage/.style={draw,rounded corners=5pt,minimum width=11.5cm,minimum height=2.4cm,align=left,inner sep=8pt},
    stitle/.style={font=\bfseries\small,anchor=north west},
    arr/.style={-{Stealth[length=3mm]},very thick,blue!50!black},
    every node/.style={font=\small}
]

% Étape 1
\node[stage,fill=red!5] (s1) at (0,0) {};
\node[stitle] at (s1.north west) {Étape 1\,: Distillation};
\node[anchor=north west,text width=10.5cm] at ($(s1.north west)+(0.1,-0.55)$) {%
\begin{tikzpicture}[
    bx/.style={draw,rounded corners=2pt,minimum height=0.7cm,align=center,font=\footnotesize},
    ar/.style={-{Stealth[length=2mm]},thick}
]
\node[bx,fill=red!15,text width=2.5cm] (teacher) {Enseignant (PINN)\\$\mathcal{L}= \text{MSE}_{\text{EDP}}+\text{MSE}_{\text{CL}}$};
\node[bx,fill=orange!12,text width=2.5cm,right=1.5cm of teacher] (student) {Élève (MLP)\\$\mathcal{L}= \text{MSE}_T + \lambda\|\theta\|_2^2$};
\node[right=0.6cm of student,font=\footnotesize,text width=3.5cm] (note) {Pas de perte EDP sur l'élève $\Rightarrow$ la rég.\ L2 \textbf{ne conflit pas} avec la physique.};
\draw[ar] (teacher) -- node[above,font=\scriptsize]{sortie $\tilde{u}_T$} (student);
\end{tikzpicture}};

% Étape 2
\node[stage,fill=yellow!5,below=0.5cm of s1] (s2) {};
\node[stitle] at (s2.north west) {Étape 2\,: Extraction de la structure};
\node[anchor=north west,text width=10.5cm] at ($(s2.north west)+(0.1,-0.55)$) {%
\begin{tikzpicture}[
    bx/.style={draw,rounded corners=2pt,minimum height=0.7cm,align=center,font=\footnotesize},
    ar/.style={-{Stealth[length=2mm]},thick}
]
\node[bx,fill=yellow!15,text width=2.2cm] (reg) {Rég.\ L2 $\Rightarrow$\\params convergent};
\node[bx,fill=yellow!20,text width=2.2cm,right=0.8cm of reg] (hac) {Clustering HAC\\sur $|\theta_l|$};
\node[bx,fill=yellow!25,text width=2.8cm,right=0.8cm of hac] (lr) {Matrices rang faible\\$\W_l = f(W^a, W^b,\ldots)$};
\draw[ar] (reg) -- (hac);
\draw[ar] (hac) -- node[above,font=\scriptsize]{remplacer par $\mu_k$} (lr);
\end{tikzpicture}};

% Étape 3
\node[stage,fill=green!5,below=0.5cm of s2] (s3) {};
\node[stitle] at (s3.north west) {Étape 3\,: Reconstruction};
\node[anchor=north west,text width=10.5cm] at ($(s3.north west)+(0.1,-0.55)$) {%
\begin{tikzpicture}[
    bx/.style={draw,rounded corners=2pt,minimum height=0.7cm,align=center,font=\footnotesize},
    ar/.style={-{Stealth[length=2mm]},thick}
]
\node[bx,fill=green!12,text width=2.5cm] (rel) {Construire la matrice\\de relation $\R_l$};
\node[bx,fill=green!18,text width=2.8cm,right=0.8cm of rel] (embed) {Contraindre le réseau\\$\hat{\theta}_l = \R_l \cdot \bar{\theta}_l$};
\node[bx,fill=green!25,text width=2.5cm,right=0.8cm of embed] (train) {Entraîner \PsiNN{}\\sur l'EDP originale};
\draw[ar] (rel) -- (embed);
\draw[ar] (embed) -- (train);
\end{tikzpicture}};

\draw[arr] (s1.south) -- (s2.north);
\draw[arr] (s2.south) -- (s3.north);

\end{tikzpicture}
\caption{Pipeline complet de \PsiNN{}.}
\end{figure}

%=================================================================
\section{Étape 1 -- Distillation\,: comment et pourquoi}
%=================================================================

\textbf{Problème central\,:} Ajouter une régularisation (L1, L2, GrOWL) à un PINN \emph{augmente} la perte car le gradient de régularisation s'oppose au gradient de l'EDP.

\textbf{Solution\,:} Séparer les deux tâches dans deux réseaux\,:

\begin{center}
\renewcommand{\arraystretch}{1.3}
\begin{tabular}{@{}l l l@{}}
\toprule
& \textbf{Enseignant} & \textbf{Élève} \\
\midrule
Architecture & MLP & MLP (même taille) \\
Perte & $\text{MSE}_{\text{EDP}} + \text{MSE}_{\text{CL}} + \text{MSE}_{\text{données}}$ & $\text{MSE}_T + \displaystyle\sum_l \omega_l\|\theta_l\|_2^2$ \\
Dérivées EDP & Oui (ordre élevé) & \textbf{Aucune} \\
Régularisation & Aucune & L2 (favorise la convergence) \\
Rôle & Apprendre la physique & Apprendre la sortie + révéler la structure \\
\bottomrule
\end{tabular}
\end{center}

\textbf{Perte de l'élève\,:}
\begin{equation}
\mathcal{L}_{\text{élève}} = \frac{1}{M_T}\sum_{i=1}^{M_T} \big|\tilde{u}_T^{(i)} - \tilde{u}_S^{(i)}\big|^2 + \sum_{l=1}^{L} \omega_l \|\theta_l\|_2^2
\end{equation}

%=================================================================
\section{Étape 2 -- Extraction\,: comment et pourquoi}
%=================================================================

\subsection{Fondement théorique}

\begin{tcolorbox}[colback=yellow!6,colframe=orange!60!black,title={\small Théorème 1\,: Les paramètres équivalents convergent}]
\small Les paramètres $\theta_1,\dots,\theta_n$ d'une même couche jouant des \textbf{rôles équivalents} $\xrightarrow{\text{rég.\ L2}}$ convergent vers la \textbf{même valeur}.
\end{tcolorbox}

\begin{tcolorbox}[colback=yellow!6,colframe=orange!60!black,title={\small Théorème 2\,: Les paramètres symétriques convergent en valeur absolue}]
\small Si $|\theta_i|=|\theta_j|$ (possiblement de signes différents) $\xrightarrow{\text{rég.\ L2}}$ leurs valeurs absolues convergent.
\end{tcolorbox}

\subsection{Étapes d'extraction}

\begin{enumerate}[leftmargin=1.8em,itemsep=2pt]
\item Pour chaque couche $l$, calculer $|\theta_l|$.
\item Exécuter le HAC (clustering hiérarchique agglomératif) $\longrightarrow$ $K$ clusters $\{C_1,\dots,C_K\}$.
\item Remplacer chaque poids par le centre de son cluster, en gardant le signe\,:
\begin{equation}
    \theta_{l,i}^{\text{nouveau}} = \operatorname{sgn}(\theta_{l,i})\cdot\mu_k, \quad \theta_{l,i}\in C_k
\end{equation}
\item Les matrices de poids montrent un \textbf{motif de rang faible}\,: chaque matrice complète est assemblée à partir de quelques sous-matrices indépendantes via partage, inversion de signe et permutations.
\end{enumerate}

\subsection{Exemple Laplace -- matrices extraites}

\begin{equation}
\W_1\!=\!\begin{bmatrix} W_1^a & W_1^b \\ -W_1^a & -W_1^b \end{bmatrix}\!,\;
\W_2\!=\!\begin{bmatrix} -W_2^a & -W_2^b \\ W_2^b & W_2^a \end{bmatrix}\!,\;
\W_3\!=\!\begin{bmatrix} W_3^a & W_3^a \end{bmatrix}
\end{equation}

\begin{center}
\begin{tikzpicture}[
    every node/.style={font=\footnotesize},
    mat/.style={draw,minimum width=1.4cm,minimum height=0.55cm,align=center,rounded corners=1pt}
]
\node[mat,fill=blue!12] (wa) {$W^a$};
\node[mat,fill=red!12,right=0.3cm of wa] (mwa) {$-W^a$};
\node[mat,fill=green!12,right=0.3cm of mwa] (wb) {$W^b$};
\node[mat,fill=purple!12,right=0.3cm of wb] (mwb) {$-W^b$};
\node[right=0.3cm of mwb] {$\leftarrow$ seulement 2 sous-matrices libres};
\end{tikzpicture}
\end{center}

%=================================================================
\section{Étape 3 -- Reconstruction\,: comment et pourquoi}
%=================================================================

\subsection{La matrice de relation $\R$}

$\R$ encode trois types de relations entre paramètres\,:

\begin{center}
\begin{tikzpicture}[
    every node/.style={font=\small},
    box/.style={draw,rounded corners=2pt,minimum height=0.6cm,minimum width=2cm,align=center}
]
\node[box,fill=blue!10]   (s) at (0,0) {$+1$\,: partage};
\node[box,fill=red!10]    (r) at (4,0) {$-1$\,: inversion};
\node[box,fill=gray!10]   (u) at (8,0) {$0$\,: non relié};
\node[below=0.1cm of s,font=\scriptsize,gray] {même paramètre};
\node[below=0.1cm of r,font=\scriptsize,gray] {signe opposé};
\node[below=0.1cm of u,font=\scriptsize,gray] {aucune connexion};
\end{tikzpicture}
\end{center}

Les permutations de lignes dans $\R$ encodent les relations d'\textbf{échange}.

\subsection{Exemple Laplace -- 2\textsuperscript{e} couche cachée}

\begin{equation}
\underbrace{\R_2 = \begin{bmatrix} -1 & 0 \\ 0 & -1 \\ 0 & 1 \\ 1 & 0 \end{bmatrix}}_{\text{encode la structure}}
\;\times\;
\underbrace{\begin{bmatrix} \hat{\theta}_2^a & 0 \\ 0 & \hat{\theta}_2^b \end{bmatrix}}_{\text{params libres (seul.\ 2)}}
\;=\;
\underbrace{\begin{bmatrix} -\hat{\theta}_2^a & 0 \\ 0 & -\hat{\theta}_2^b \\ 0 & \hat{\theta}_2^b \\ \hat{\theta}_2^a & 0 \end{bmatrix}}_{\text{matrice complète contrainte}}
\end{equation}

\subsection{Procédure de reconstruction}

\begin{center}
\begin{tikzpicture}[
    step/.style={draw,rounded corners=3pt,fill=green!8,minimum height=0.7cm,text width=3.2cm,align=center,font=\footnotesize},
    arr/.style={-{Stealth[length=2mm]},thick}
]
\node[step] (a) at (0,0) {Sparsifier\\(mettre le bruit à zéro)};
\node[step,right=0.5cm of a] (b) {Réinitialiser\\les params saillants};
\node[step,right=0.5cm of b] (c) {Appliquer $\R$\\(contraindre les poids)};
\node[step,right=0.5cm of c] (d) {Entraîner \PsiNN{}\\sur l'EDP originale};
\draw[arr](a)--(b); \draw[arr](b)--(c); \draw[arr](c)--(d);
\end{tikzpicture}
\end{center}

%=================================================================
\section{Pourquoi ça marche\,: preuve de symétrie (Laplace)}
%=================================================================

Avec la structure extraite, les deux sorties de la dernière couche cachée sont\,:
\begin{align}
u^a(x_1,x_2) &= \tanh\!\big(W_2^a\,\sigma_+ + W_2^b\,(\sigma_- + b_2^a)\big) \\
u^b(x_1,x_2) &= \tanh\!\big(W_2^b\,\sigma_+ + W_2^a\,(\sigma_- + b_2^a)\big)
\end{align}
où $\sigma_\pm = \tanh(W_1^a x_1 \pm W_1^b x_2 + b_1^a)$.

Remplacer $x_2\to -x_2$ échange $\sigma_+\leftrightarrow\sigma_-$, donc\,:
\begin{equation}
\boxed{u^a(x_1,x_2) = u^b(x_1,-x_2)}
\end{equation}

Sortie finale\,:
\begin{equation}
u(x_1,x_2) = W_3^a\big(u^a(x_1,x_2)+u^a(x_1,-x_2)\big) \;\Longrightarrow\; u(x_1,x_2)=u(x_1,-x_2)
\end{equation}

Le réseau \textbf{satisfait intrinsèquement} la symétrie de Laplace -- aucun post-traitement nécessaire.

%=================================================================
\section{Résultats numériques}
%=================================================================

\subsection{Quatre cas de référence}

\begin{center}
\renewcommand{\arraystretch}{1.25}
\begin{tabular}{@{}l l l l@{}}
\toprule
\textbf{Cas} & \textbf{EDP} & \textbf{Symétrie trouvée} & \textbf{Particularité} \\
\midrule
Laplace    & $\nabla^2 u=0$ & $u(x_1,x_2)=u(x_1,-x_2)$ & Walkthrough complet d'extraction \\
Burgers    & $u_t+uu_x+\lambda_1 u_{xx}=0$ & Antisymétrique & Prob.\ inverse\,; transfère entre $\nu$ \\
Poisson    & $\nabla^2 u=f$ (4 fréquences) & Équivariance $\mathbb{Z}_2\!\times\!\mathbb{Z}_2$ & Transfert basse$\to$haute fréquence \\
Cylindre   & Navier--Stokes & Symétries opposées sur $u,v,p$ & Multi-sortie\,; réutilise les structures \\
\bottomrule
\end{tabular}
\end{center}

\subsection{Comparaison des erreurs $L_2$}

\begin{center}
\renewcommand{\arraystretch}{1.2}
\begin{tabular}{@{}l c c c@{}}
\toprule
\textbf{Problème} & \textbf{PINN} & \textbf{PINN-post} & \textbf{\PsiNN{}} \\
\midrule
Laplace ($\times10^{-4}$)       & 11,59 & 4,017 & \textbf{0,742} \\
Burgers ($\times10^{-2}$)       & 14,47 & 3,014 & \textbf{1,287} \\
Poisson ($\times10^{-2}$)       & 2,633 & 2,563 & \textbf{2,464} \\
Écoul.\ $p$ ($\times10^{-4}$)   & 14,89 & 11,78 & \textbf{7,838} \\
Écoul.\ $u$ ($\times10^{-4}$)   & 1,981 & 1,904 & \textbf{1,854} \\
Écoul.\ $v$ ($\times10^{-5}$)   & 1,984 & 1,896 & \textbf{1,765} \\
\bottomrule
\end{tabular}
\end{center}

\PsiNN{} obtient l'\textbf{erreur la plus faible dans chaque cas}.

\subsection{Burgers -- estimation du paramètre inverse}

\begin{center}
\renewcommand{\arraystretch}{1.2}
\begin{tabular}{@{}l c c c c@{}}
\toprule
& \textbf{Vérité} & \textbf{PINN} & \textbf{PINN-post} & \textbf{\PsiNN{}} \\
\midrule
$\nu_1\!\cdot\!\pi$ ($\times10^{-2}$) & 1 & 2,820 & 1,898 & \textbf{1,465} \\
$\nu_2\!\cdot\!\pi$ ($\times10^{-2}$) & 4 & 6,625 & 4,960 & \textbf{4,084} \\
$\nu_3\!\cdot\!\pi$ ($\times10^{-2}$) & 8 & 9,705 & 0,885 & \textbf{6,673} \\
\bottomrule
\end{tabular}
\end{center}

Même structure \PsiNN{} utilisée pour les trois valeurs de $\nu$ -- aucun ré-entraînement nécessaire.

%=================================================================
\section{Avantages vs.\ limites}
%=================================================================

\begin{center}
\renewcommand{\arraystretch}{1.3}
\begin{tabular}{@{}p{6.5cm} p{6.5cm}@{}}
\toprule
\textbf{Avantages} & \textbf{Limites} \\
\midrule
Découverte automatique (pas de connaissances a priori) & Nécessite des EDP au moins partiellement connues \\
Physique \textit{dans} les poids $\Rightarrow$ satisfaction stricte & Validé uniquement sur MLP 3 couches + $\tanh$ \\
Interprétable (matrices rang faible $\leftrightarrow$ symétries) & Limité aux symétries de signe \& permutation \\
Moins de paramètres $\Rightarrow$ entraînement plus rapide & Extension aux Transformers = travaux futurs \\
Transférable entre paramètres \& problèmes & Symétries complexes (rotation, Noether) non traitées \\
\bottomrule
\end{tabular}
\end{center}

%=================================================================
\section{Équations clés -- Référence rapide}
%=================================================================

\begin{tcolorbox}[colback=blue!3,colframe=blue!40!black]
\renewcommand{\arraystretch}{1.8}
\begin{tabular}{@{}l l@{}}
\textbf{Perte de distillation} & $\displaystyle\text{MSE}_T = \frac{1}{M_T}\sum_{i=1}^{M_T}|\tilde{u}_T^{(i)}-\tilde{u}_S^{(i)}|^2$ \\
\textbf{Régularisation L2} & $\displaystyle\Omega(\theta_S) = \sum_{l=1}^{L}\omega_l\|\theta_l\|_2^2$ \\
\textbf{Remplacement par clustering} & $\theta_{l,i}^{\text{nouveau}} = \operatorname{sgn}(\theta_{l,i})\cdot\mu_k \;\;(\theta_{l,i}\in C_k)$ \\
\textbf{Reconstruction} & $\hat{\theta}_l = \R_l\cdot\bar{\theta}_l \;\;$ où $\R_l\in\{-1,0,+1\}^{n\times K}$ \\
\end{tabular}
\end{tcolorbox}

\end{document}
